\documentclass[11pt,a4paper]{article}

\usepackage[utf8]{inputenc}
\usepackage[T1]{fontenc}
\usepackage[french]{babel}
\usepackage[a4paper,margin=2.5cm]{geometry}
\usepackage{parskip}
\usepackage{hyperref}
\usepackage{enumitem}
\usepackage{graphicx}
\usepackage{float}
\usepackage{xcolor}
\usepackage[normalem]{ulem}
\usepackage{amsmath}
\usepackage{booktabs}
\usepackage{array}

\title{Rapport de Progrès 18/05/2026}
\author{Xihao WANG}
\date{\today}

\begin{document}

\maketitle

\section{Contexte et rappel bref}

L'objectif du projet reste le suivi multi-objets dans une même caméra lorsque l'apparence d'une personne change fortement au cours du temps. Le problème étudié n'est donc pas seulement de détecter correctement les personnes, mais de conserver une identité cohérente malgré les changements de vêtement, les occlusions, les sorties temporaires du champ et les ambiguïtés entre plusieurs individus proches.

À la fin du rapport précédent, le travail avait permis :
\begin{itemize}
    \item d'établir une baseline fondée sur \textit{DeepSORT + BoT} ;
    \item d'ajouter des mécanismes de mémoire \textit{short-term} et \textit{long-term} ;
    \item de remplacer progressivement un \textit{temporal trend} manuel par un module temporel appris ;
    \item de construire un premier \textit{debug viewer} pour inspecter les matrices de coût et les scores temporels image par image.
\end{itemize}

Le travail de ce mois-ci a porté moins sur la simple preuve de concept du module temporel que sur son intégration dans un tracker complet, la fiabilisation des données d'entraînement et l'évaluation globale du système avec des métriques standardisées.

\section{Travail réalisé ce mois-ci}

\subsection{Passage d'une logique inspirée de TIE vers une structure de type mémoire}

La première tentative de module appris suivait encore une logique proche de \textit{TIE} : le modèle cherchait surtout à décrire l'évolution locale de l'apparence à partir de quelques états temporels récents. Cette piste a permis de vérifier que l'information temporelle pouvait être exploitée, mais elle s'est révélée trop limitée dans notre cadre multi-objets : les tokens proches restaient souvent très similaires entre eux, l'attention apportait peu de séparation utile, et le modèle généralisait mal.

Le travail de ce mois-ci a donc marqué un changement de formulation. Au lieu de poursuivre une architecture centrée sur quelques différences temporelles locales, le modèle a été reconstruit dans une logique plus proche de \textit{MeMOT}, avec :
\begin{itemize}
    \item une mémoire courte pour représenter l'apparence récente ;
    \item une mémoire longue pour conserver une information d'identité plus stable ;
    \item un score appris qui compare directement la détection courante avec ces deux mémoires.
\end{itemize}

Cette transition explique l'architecture actuelle du système : le module appris n'est plus seulement un estimateur de tendance locale, mais un composant de matching fondé sur des mémoires temporelles.

\subsection{Construction d'un pipeline de données plus fiable}

La première amélioration importante a été de remplacer les pseudo-labels issus des sorties du tracker par des annotations \textit{ground truth} préparées manuellement. Cette décision était nécessaire, car une erreur du tracker utilisée comme pseudo-label peut être réinjectée dans l'entraînement et apprendre au modèle une association incorrecte.

Pour cela, plusieurs outils ont été ajoutés :
\begin{itemize}
    \item un mode d'édition de \textit{ground truth} dans le \textit{viewer}, permettant de déplacer les boîtes, modifier les identifiants, fusionner des identités et exporter directement un fichier \texttt{gt.txt} au format MOT ;
    \item un générateur de paires temporelles fondé sur le \textit{ground truth}, qui associe les détections aux identités réelles par IoU et conserve une mémoire par identité même après une absence prolongée ;
\end{itemize}

Le dataset d'apprentissage contient les volumes suivants :

\begin{table}[H]
\centering
\begin{tabular}{l r r r}
\toprule
\textbf{Vidéo} & \textbf{Paires totales} & \textbf{Positives} & \textbf{Négatives} \\
\midrule
new-1 & 2344 & 1263 & 1081 \\
new-2 & 2819 & 985 & 1834 \\
new-3 & 2100 & 1185 & 915 \\
new-5 & 2526 & 1296 & 1230 \\
new-6 & 2807 & 1147 & 1660 \\
new-8 & 3348 & 1674 & 1674 \\
new-9 & 1340 & 750 & 590 \\
\midrule
YT-03 (validation) & 5180 & 2625 & 2555 \\
\bottomrule
\end{tabular}
\caption{Répartition des paires temporelles construites à partir du \textit{ground truth}.}
\end{table}

Cette construction permet désormais d'apprendre directement à partir d'identités réelles et non plus à partir des identifiants parfois fragmentés de la baseline.

\subsection{Intégration du score appris dans le matching réel}

La baseline produit une distance d'apparence , alors que le modèle appris produit une probabilité. 

Le score du modèle appris est ensuite combiné avec la distance d'apparence de la baseline, tout en tenant compte du \textit{Kalman gating}, qui apporte l'information de mouvement. Le système produit ainsi une distance finale : plus cette distance est faible, plus l'association entre un track et une détection est favorisée.

La distance de la baseline reste la base du calcul, et le modèle appris intervient comme une correction additive contrôlée. Le coefficient de fusion associé au modèle appris est fixé à \(\gamma = 1.0\) ; il ne s'agit donc pas d'une moyenne directe de type \(50\% / 50\%\), mais d'une distance baseline ajustée par le signal temporel appris.

\subsection{Gestion des sorties du champ et réactivation d'identité}

Une autre partie du travail a porté sur le cas où une personne quitte le champ de la caméra puis revient plus tard. Au lieu de supprimer complètement son historique, le système conserve désormais les informations associées à son ancien track. Lorsqu'une nouvelle personne apparaît ensuite dans l'image, elle peut donc non seulement être comparée aux tracks actuellement actifs ou activer un nouveau track, mais aussi aux tracks historiques déjà observés auparavant.

Cette extension permet d'étudier plus directement la conservation d'identité après une absence prolongée, ce qui est particulièrement important dans le cadre du changement d'apparence.


Enfin, \textit{TrackEval} a été intégré dans le dépôt afin de produire des métriques de tracking comparables et reproductibles sur plusieurs vidéos annotées.

\section{Architecture actuelle du système}

\subsection{Vue d'ensemble du tracker}

L'architecture du modèle complet peut être résumée comme suit :

\begin{figure}[H]
    \centering
    \includegraphics[width=0.95\linewidth]{1.png}
    \caption{Architecture générale du système de tracking après l'étape de détection.}
\end{figure}

Plus précisément :
\begin{itemize}
    \item à l'instant \(t-1\), chaque track possède déjà une mémoire d'apparence. Celle-ci est séparée en une mémoire courte, qui décrit l'apparence récente, et une mémoire longue, qui conserve une information plus stable sur l'identité ;
    \item à l'instant \(t\), les nouvelles détections et leurs caractéristiques visuelles constituent la deuxième entrée du système ;
    \item la branche de gauche correspond à la baseline : les mémoires courte et longue sont utilisées pour calculer une distance d'apparence \textit{memory-aware} entre chaque track et chaque détection ;
    \item en parallèle, la branche centrale correspond au module temporel appris. Il reçoit les tokens issus des deux mémoires ainsi que la caractéristique courante de la détection, puis produit un score temporel de compatibilité ;
    \item la branche de droite correspond à l'information de mouvement : la position des boîtes et l'état dynamique sont traités par le \textit{Kalman motion gating} ;
    \item ces trois informations sont ensuite réunies dans une distance finale fusionnée. Plus cette distance est faible, plus la paire \((track, detection)\) est considérée comme probable ;
    \item l'association finale est obtenue par l'algorithme hongrois, puis les tracks et leurs mémoires sont mis à jour pour l'image suivante.
\end{itemize}

\vspace{2em}

\begin{table}[H]
\centering
\begin{tabular}{l c}
\toprule
\textbf{Paramètre de la baseline} & \textbf{Valeur utilisée} \\
\midrule
\texttt{max\_age} & 100 \\
\texttt{short\_history\_len} & 5 \\
\texttt{long\_history\_len} & 15 \\
\bottomrule
\end{tabular}
\caption{Principaux paramètres utilisés pour la baseline et les mémoires du tracker.}
\end{table}

\begin{table}[H]
\centering
\begin{tabular}{l c}
\toprule
\textbf{Paramètre du modèle temporel appris} & \textbf{Valeur utilisée} \\
\midrule
\texttt{hidden\_dim} & 256 \\
\texttt{num\_heads} & 4 \\
\texttt{batch\_size} & 64 \\
\texttt{learning\_rate} & $1 \times 10^{-4}$ \\
\texttt{epochs} & 50 \\
\bottomrule
\end{tabular}
\caption{Principaux paramètres utilisés pour l'entraînement du modèle temporel appris.}
\end{table}

\subsection{Structure du modèle temporel appris}

La structure interne du module temporel appris est résumée dans la figure suivante :

\begin{figure}[H]
    \centering
    \includegraphics[width=1.0\linewidth]{2.png}
    \caption{Architecture du module temporel appris utilisé dans la configuration finale.}
\end{figure}


Le modèle ne repose pas seulement sur les vecteurs encodés par attention. La tête de fusion reçoit aussi plusieurs raccourcis explicites, notamment :
\begin{itemize}
    \item la similarité entre la détection et le résumé court ;
    \item la similarité entre la détection et le résumé long ;
    \item la valeur de la porte longue ;
    \item des similarités cosinus brutes calculées directement entre la détection et les tokens de mémoire courte.
\end{itemize}

Ce choix s'est révélé important en pratique : dans plusieurs cas, les poids d'attention seuls ne suffisent pas à expliquer la séparation finale entre deux candidats, alors que les similarités explicites et la tête de fusion portent une partie substantielle du signal discriminant.



\subsection{Objectif d'entraînement}

Le modèle est entraîné sur des paires \((track, detection)\). Chaque paire reçoit :
\[
label =
\begin{cases}
    1, & \text{si la détection appartient à la même identité réelle,} \\
    0, & \text{sinon.}
\end{cases}
\]
L'objectif final combine deux termes :
\begin{itemize}
    \item une perte binaire \texttt{BCEWithLogitsLoss} ;
    \item une perte de classement entre candidats d'une même détection, imposant que le bon track reçoive un score supérieur aux compétiteurs négatifs difficiles.
\end{itemize}

La perte de ranking impose schématiquement :
\[
score(track^+, detection) \geq score(track^-, detection) + margin
\]
avec une marge de $0.2$ et une sélection du négatif le plus difficile par groupe.

L'entraînement final sur les paires GT a obtenu :
\begin{itemize}
    \item meilleure perte de validation : \textbf{0.0392} ;
    \item époque correspondante : \textbf{46} ;
    \item meilleure précision de validation observée : \textbf{0.9894}.
\end{itemize}

\section{Évaluation expérimentale}

\subsection{Métriques utilisées}

L'évaluation globale du tracker a été réalisée avec \textit{TrackEval}. Les principales métriques retenues sont :
\begin{itemize}
    \item \textbf{HOTA} : mesure globale combinant détection et association ;
    \item \textbf{DetA} : qualité de la détection ;
    \item \textbf{AssA} : qualité de l'association des identités au cours du suivi ;
    \item \textbf{DetRe / DetPr} : rappel et précision de détection ;
    \item \textbf{AssRe / AssPr} : rappel et précision d'association ;
    \item \textbf{LocA} : précision de localisation spatiale ;
    \item \textbf{IDF1} : cohérence globale des identités ;
    \item \textbf{IDR / IDP} : rappel et précision au niveau des identités ;
    \item \textbf{MOTA} : métrique CLEAR MOT globale ;
    \item \textbf{IDSW} : nombre de changements d'identité ;
    \item \textbf{FP / FN} : faux positifs et faux négatifs.
\end{itemize}

Deux configurations ont été comparées sur six vidéos annotées :
\begin{itemize}
    \item \textbf{sans réactivation} : fusion du modèle temporel avec le tracker, sans réactivation ;
    \item \textbf{avec réactivation} : même configuration, avec réactivation inactive, score appris dans le gate et probation après réactivation.
\end{itemize}

\subsection{Résultats principaux}

Les résultats présentés ci-dessous correspondent aux sorties du modèle complet, c'est-à-dire après la fusion entre la baseline, le module temporel appris et les contraintes de mouvement.

\begin{table}[H]
\centering
\resizebox{\textwidth}{!}{
\begin{tabular}{l l r r r r r r r r}
\toprule
\textbf{Vidéo} & \textbf{Configuration} & \textbf{HOTA} & \textbf{DetA} & \textbf{AssA} & \textbf{DetRe} & \textbf{DetPr} & \textbf{AssRe} & \textbf{AssPr} & \textbf{LocA} \\
\midrule
YT-03 & sans réactivation & 90.557 & 96.052 & 85.379 & 99.492 & 96.509 & 85.914 & 96.941 & 99.459 \\
YT-03 & avec réactivation & \textbf{91.470} & 95.384 & \textbf{87.718} & 99.492 & 95.836 & \textbf{88.830} & 96.448 & 99.458 \\
\midrule
YT-08 & sans réactivation & \textbf{89.811} & \textbf{98.700} & \textbf{81.723} & 99.344 & \textbf{99.344} & 87.605 & \textbf{88.911} & 99.752 \\
YT-08 & avec réactivation & 89.548 & 98.437 & 81.461 & 99.344 & 99.077 & 87.605 & 88.649 & 99.752 \\
\midrule
YT-09 & sans réactivation & 99.655 & 99.652 & 99.658 & 99.824 & 99.824 & 99.826 & 99.826 & 99.732 \\
YT-09 & avec réactivation & 99.655 & 99.652 & 99.658 & 99.824 & 99.824 & 99.826 & 99.826 & 99.732 \\
\midrule
new-4 & sans réactivation & \textbf{66.072} & \textbf{99.819} & 43.734 & \textbf{99.939} & \textbf{99.879} & \textbf{49.893} & 84.031 & \textbf{99.958} \\
new-4 & avec réactivation & 65.999 & 99.491 & \textbf{43.782} & 99.805 & 99.685 & 49.822 & \textbf{84.046} & 99.907 \\
\midrule
new-7 & sans réactivation & 90.937 & 90.706 & 91.172 & 97.913 & 92.333 & 97.200 & 93.219 & 98.070 \\
new-7 & avec réactivation & 90.937 & 90.706 & 91.172 & 97.913 & 92.333 & 97.200 & 93.219 & 98.070 \\
\midrule
new-10 & sans réactivation & 99.241 & 99.309 & 99.174 & 99.647 & 99.647 & 99.549 & 99.551 & 99.566 \\
new-10 & avec réactivation & 99.241 & 99.309 & 99.174 & 99.647 & 99.647 & 99.549 & 99.551 & 99.566 \\
\bottomrule
\end{tabular}
}
\caption{Résultats HOTA détaillés avec et sans réactivation.}
\end{table}

\begin{table}[H]
\centering
\resizebox{\textwidth}{!}{
\begin{tabular}{l l r r r r r r r r r}
\toprule
\textbf{Vidéo} & \textbf{Configuration} & \textbf{MOTA} & \textbf{IDF1} & \textbf{IDR} & \textbf{IDP} & \textbf{FP} & \textbf{FN} & \textbf{IDSW} & \textbf{Frag} & \textbf{IDs prédits} \\
\midrule
YT-03 & sans réactivation & \textbf{96.566} & 85.704 & 87.028 & 84.419 & 83 & 2 & 5 & 20 & 7 \\
YT-03 & avec réactivation & 95.918 & \textbf{90.303} & \textbf{92.026} & \textbf{88.644} & 102 & 2 & \textbf{3} & 20 & \textbf{5} \\
\midrule
YT-08 & sans réactivation & \textbf{99.892} & \textbf{91.559} & 91.559 & \textbf{91.559} & \textbf{0} & 0 & 2 & 26 & 4 \\
YT-08 & avec réactivation & 99.624 & 91.436 & 91.559 & 91.314 & 5 & 0 & 2 & 26 & 4 \\
\midrule
YT-09 & sans réactivation & 100.000 & 100.000 & 100.000 & 100.000 & 0 & 0 & 0 & 18 & 4 \\
YT-09 & avec réactivation & 100.000 & 100.000 & 100.000 & 100.000 & 0 & 0 & 0 & 18 & 4 \\
\midrule
new-4 & sans réactivation & \textbf{99.758} & \textbf{56.104} & 56.121 & \textbf{56.087} & \textbf{1} & \textbf{0} & 3 & 19 & 4 \\
new-4 & avec réactivation & 99.455 & 56.087 & 56.121 & 56.053 & 4 & 2 & 3 & 19 & 4 \\
\midrule
new-7 & sans réactivation & 92.715 & 96.183 & 99.089 & 93.443 & 80 & 7 & 1 & 19 & 2 \\
new-7 & avec réactivation & 92.715 & 96.183 & 99.089 & 93.443 & 80 & 7 & 1 & 19 & 2 \\
\midrule
new-10 & sans réactivation & 99.823 & 99.912 & 99.912 & 99.912 & 0 & 0 & 2 & 14 & 2 \\
new-10 & avec réactivation & 99.823 & 99.912 & 99.912 & 99.912 & 0 & 0 & 2 & 14 & 2 \\
\bottomrule
\end{tabular}
}
\caption{Résultats CLEAR MOT et Identity avec et sans réactivation.}
\end{table}

\begin{table}[H]
\centering
\begin{tabular}{l r r r r}
\toprule
\textbf{Configuration} & \textbf{HOTA moyen} & \textbf{AssA moyen} & \textbf{IDF1 moyen} & \textbf{IDSW total} \\
\midrule
sans réactivation & 89.379 & 83.473 & 88.244 & 13 \\
avec réactivation & \textbf{89.475} & \textbf{83.828} & \textbf{88.987} & \textbf{11} \\
\bottomrule
\end{tabular}
\caption{Moyenne non pondérée sur les six vidéos évaluées.}
\end{table}

\subsection{Évaluation ciblée sur les identités concernées par le changement d'apparence}

Afin d'analyser plus directement le cas du changement d'apparence, une seconde évaluation a été réalisée sur des identités ciblées, en évitant autant que possible les personnes dont les changements d'ID proviennent surtout d'occlusions ou de sorties du champ de la caméra.

Les identités évaluées sont :
\begin{itemize}
    \item \textit{GT1} pour \textit{new-4}, \textit{new-7}, \textit{YT-03} et \textit{new-10} ;
    \item \textit{GT2} pour \textit{YT-08} ;
    \item \textit{GT1}, \textit{GT2} et \textit{GT3} pour \textit{YT-09}.
\end{itemize}

\begin{table}[H]
\centering
\resizebox{\textwidth}{!}{
\begin{tabular}{l l r r r r}
\toprule
\textbf{Vidéo / cible} & \textbf{Configuration} & \textbf{HOTA} & \textbf{AssA} & \textbf{IDF1} & \textbf{IDSW} \\
\midrule
new-4 / GT1 & sans réactivation & 71.057 & 50.490 & 54.952 & 1 \\
new-4 / GT1 & avec réactivation & 71.057 & 50.490 & 54.952 & 1 \\
\midrule
new-7 / GT1 & sans réactivation & 96.138 & 95.800 & 99.145 & 3 \\
new-7 / GT1 & avec réactivation & 96.138 & 95.800 & 99.145 & 3 \\
\midrule
YT-03 / GT1 & sans réactivation & 99.126 & 99.126 & 99.954 & 0 \\
YT-03 / GT1 & avec réactivation & 99.126 & 99.126 & 99.954 & 0 \\
\midrule
new-10 / GT1 & sans réactivation & 99.456 & 99.317 & 99.771 & 4 \\
new-10 / GT1 & avec réactivation & 99.456 & 99.317 & 99.771 & 4 \\
\midrule
YT-08 / GT2 & sans réactivation & 100.000 & 100.000 & 100.000 & 0 \\
YT-08 / GT2 & avec réactivation & 100.000 & 100.000 & 100.000 & 0 \\
\midrule
YT-09 / GT1+GT2+GT3 & sans réactivation & 99.575 & 99.578 & 100.000 & 0 \\
YT-09 / GT1+GT2+GT3 & avec réactivation & 99.575 & 99.578 & 100.000 & 0 \\
\bottomrule
\end{tabular}
}
\caption{Évaluation TrackEval ciblée sur les identités d'intérêt pour l'analyse du changement d'apparence.}
\end{table}

\subsection{Précision propre du modèle temporel appris}

En complément du tracker complet, le modèle temporel appris a aussi été évalué isolément à partir des matrices de scores reconstruites sur les données \textit{ground truth}. Pour chaque frame, seuls les scores du modèle appris sont utilisés. La métric correspond au taux de précision où toutes les associations cibles sont correctement identifiées.

\begin{table}[H]
\centering
\begin{tabular}{l r}
\toprule
\textbf{Vidéo / cible} & \textbf{Frame Hungarian accuracy} \\
\midrule
new-4 / GT1 & 100.000 \\
new-7 / GT1 & 100.000 \\
YT-03 / GT1 & 100.000 \\
new-10 / GT1 & 99.883 \\
YT-08 / GT1+GT2+GT3 & 99.855 \\
YT-09 / GT1+GT2+GT3 & 100.000 \\
\bottomrule
\end{tabular}
\caption{Précision du modèle temporel appris évalué seul par appariement hongrois sur les identités ciblées.}
\end{table}

\subsection{Analyse des résultats}

\begin{itemize}
    \item Sur les six vidéos complètes, le système obtient de très bons résultats sur \textit{YT-09} et \textit{new-10}. À l'inverse, \textit{new-4} reste le cas le plus difficile : la détection y est presque parfaite, mais l'association reste faible (\textit{AssA} autour de $43.7$, \textit{IDF1} autour de $56.1$).

    \item La réactivation apporte un gain global modéré : l'\textit{IDF1} moyen passe de $88.244$ à $88.987$, l'\textit{AssA} moyen de $83.473$ à $83.828$, et les \textit{ID switches} totaux diminuent de $13$ à $11$. Le gain vient surtout de \textit{YT-03}, tandis que la réactivation reste neutre sur \textit{YT-09}, \textit{new-7} et \textit{new-10}, et légèrement défavorable sur \textit{YT-08} et \textit{new-4}.

    \item L'évaluation ciblée montre que plusieurs identités directement concernées par le changement d'apparence sont déjà très bien suivies : \textit{YT-03/GT1}, \textit{YT-08/GT2}, \textit{YT-09/GT1+GT2+GT3} et \textit{new-10/GT1} sont proches de la perfection. Le principal échec ciblé reste \textit{new-4/GT1}, avec un \textit{IDF1} de $54.952$.

    \item Évalué seul avec un appariement hongrois fondé uniquement sur ses scores, le modèle temporel appris atteint une précision par frame presque parfaite : $100\%$ sur quatre cas, $99.883\%$ sur \textit{new-10/GT1} et $99.855\%$ sur \textit{YT-08/GT1+GT2+GT3}. Le score appris est donc déjà très discriminant ; les erreurs restantes du tracker complet viennent surtout de son intégration avec le reste du pipeline.
\end{itemize}

\section{Limites actuelles}

Malgré les progrès obtenus, plusieurs limites restent présentes :
\begin{itemize}
    \item le nombre de vidéos annotées reste encore limité, ce qui rend les conclusions statistiques prudentes ;
    \item la réactivation reste utile surtout dans les scènes comportant de vraies sorties et réapparitions, mais peut encore introduire quelques faux positifs supplémentaires ;
    \item certaines erreurs proviennent encore du couple \textit{appearance + motion gating}, notamment lorsque plusieurs candidats proches reçoivent des coûts presque équivalents.
\end{itemize}

\section{Conclusion et perspectives}

Le travail de ce mois-ci a permis de franchir une étape importante : après les limites observées avec une première approche inspirée de \textit{TIE}, le projet a évolué vers une architecture davantage fondée sur les mémoires courte et longue, plus adaptée au matching multi-objets étudié ici. Le module temporel appris n'est désormais plus seulement un score auxiliaire observé après coup, mais une composante intégrée au tracker réel. En parallèle, la création d'un pipeline d'annotation GT, d'un constructeur de paires fondé sur le GT et d'une évaluation TrackEval complète a rendu les expériences beaucoup plus fiables.

Les résultats montrent que :
\begin{itemize}
    \item le passage d'une logique de tendance temporelle locale vers une structure fondée sur la mémoire a constitué l'évolution architecturale principale de ce mois-ci ;
    \item le modèle temporel appris peut être utilisé concrètement dans la décision de matching ;
    \item l'entraînement sur des identités GT améliore fortement la qualité du signal appris ;
    \item la réactivation contrôlée est prometteuse pour les cas de sortie puis retour dans le champ, en particulier sur les métriques d'identité.
\end{itemize}

Les prochaines étapes les plus pertinentes sont :
\begin{itemize}
    \item augmenter le nombre de vidéos annotées et diversifier davantage les scénarios ;
    \item mieux traiter les réapparitions après veto, par exemple avec un état d'attente explicite avant de créer un nouveau track ;
    \item utiliser plus directement le score appris pendant la probation de réactivation ;
    \item étudier plus systématiquement les cas où le \textit{gating} de mouvement et le score temporel entrent en contradiction ;
    \item comparer quantitativement les différentes variantes de mémoire courte et longue sur un ensemble plus large.
\end{itemize}

\end{document}
